\documentclass[border=8pt]{standalone}
\usepackage[american]{circuitikz}
\begin{document}
\begin{circuitikz}[font=\sffamily]
  \node[npn, scale=1.1] (Q) at (2.3,1.8) {};
  \draw (0,0) -- (8.2,0) node[ground]{};
  \draw (0,5.0) -- (8.2,5.0) node[right]{$+V_{CC}$};
  \draw (Q.C) to[R,l=$R_C$] (2.3,5.0);
  \draw (Q.E) to[R,l_=$R_E$] (2.3,0);
  \draw (Q.B) -- (1.0,1.8) to[R,l_=$R_2$] (1.0,0);
  \draw (1.0,1.8) to[R,l=$R_1$] (1.0,5.0);
  \draw (Q.C) -- (4.0,3.0) to[C,l=$C_C$] (5.0,3.0);
  \draw (5.0,3.0) to[L,l=$L_1$] (5.0,1.5) to[L,l=$L_2$] (5.0,0);
  \draw (5.0,3.0) -- (7.0,3.0) to[C,l=$C$] (7.0,0) -- (5.0,0);
  \draw (5.0,1.5) to[C,l=$C_f$] (3.5,1.5) -- (Q.E);
  \draw (7.0,3.0) -- (8.2,3.0) node[right]{$V_o$ to CRO};
  \node[above] at (4.1,5.35) {Transistor Hartley oscillator with tapped-coil feedback};
  \node[right] at (5.1,1.5) {tap};
\end{circuitikz}
\end{document}
